\documentclass[11pt,reqno]{amsart}

\usepackage{amsmath,amssymb,amsthm}
\usepackage{booktabs}
\usepackage[margin=1in]{geometry}
\usepackage[colorlinks=true,linkcolor=blue,citecolor=blue,urlcolor=blue]{hyperref}

\setlength{\emergencystretch}{3em}

% ---------------------------------------------------------------------------
%  Perfect 1-factorisations of K_18 and K_20 with prescribed automorphisms
%  LaTeX (amsart) source. Draft version 0.6 (2026-06-30).
%
%  Changelog (not printed):
%   v0.8 (2026-07-01) Corrected Section 5.6: the two atomic order-17 squares are NOT new --
%        they are exactly the two even-starter P1Fs of K18 of Bryant-Maenhaut-Wanless (2006),
%        which our census independently rediscovers. Confirmed by reconstructing both even
%        starters E_A,E_B for p=17 and matching by paratopy (nu=27744 and nu=96). Section 5.6
%        reframed as cross-validation; abstract/contributions updated; Appendix A.6 records
%        nu and the paratopy match; ref Bryant-Maenhaut-Wanless added.
%   v0.7 (2026-07-01) Added Section 5.6 (atomic Latin squares): exactly two of the 531
%        K18 P1Fs are atomic -> two DISTINCT main classes of order-17 atomic squares,
%        vs none of order 15 for K16. Appendix A.6; refs Maenhaut-Wanless, Wanless-cyclotomic.
%   v0.6 (2026-06-30) Completed the K18 order->=3 classification: sufficient
%        sweep L in {3,4,5,7,11,13,17} finished and deduplicated => exactly 531
%        classes {3:351,4:144,8:22,16:12,17:1,272:1}. Section 5.1 rewritten
%        around this census; per-order counts (with the order-8 correction) moved
%        to a footnote. Added Section 3.5 (which orders suffice). Abstract + status
%        updated.
%   v0.5 Resolved K18 2^9 (fixed-point-free) = NON-EMPTY (Appendix A.5).
%   v0.4 Added Section 5.5 (1-rotational / starter subfamily).
%   v0.3 Added the order-2 landscape (Section 5.4).
%   v0.2 Added the K14 validation (Section 4.2).
% ---------------------------------------------------------------------------

\theoremstyle{plain}
\newtheorem{theorem}{Theorem}
\newtheorem{lemma}[theorem]{Lemma}
\newtheorem{corollary}[theorem]{Corollary}
\theoremstyle{definition}
\newtheorem*{observation}{Observation}
\newtheorem*{principle}{Automorphism-first enumeration}
\newtheorem*{ogrule}{Orderly-generation rule}

\numberwithin{theorem}{section}

\newcommand{\Aut}{\operatorname{Aut}}
\newcommand{\AGL}{\operatorname{AGL}}
\newcommand{\GF}{\operatorname{GF}}

\begin{document}

\title[Perfect 1-factorisations of $K_{18}$ and $K_{20}$]
{Perfect 1-factorisations of $K_{18}$ and $K_{20}$\\ with prescribed automorphisms}

\author{Andrei V. Ivanov}
\address{Andrei V. Ivanov, Independent Researcher, San Jose, CA, USA}

\author{Leonid M. Tertitski}
\address{Leonid M. Tertitski, Independent Researcher, Carlsbad, CA, USA}

\date{Draft --- version 0.8 (2026-07-01)}

\begin{abstract}
A \emph{perfect 1-factorisation} (P1F) of $K_{2n}$ partitions its edges into $2n-1$
perfect matchings, any two of which union to a Hamiltonian cycle. Complete enumeration
is feasible only up to $K_{16}$ (Gill and Wanless: $3155$ P1Fs, $89$ with non-trivial
automorphism group); for $K_{18}$ and beyond only isolated algebraic examples were known.

We take an \emph{automorphism-first} approach: for a fixed vertex permutation $\alpha$
we enumerate \textbf{all} $\alpha$-invariant P1Fs directly, by an exact cover of the edges
with $\alpha$-invariant orbits of 1-factors and orderly generation for isomorph rejection.
One representative $\alpha$ per cycle type of a target order, deduplicated across types,
yields every P1F with an automorphism of that order --- scaling past the range of complete
enumeration.

We validate the method by reproducing the $K_{16}$ classification of Gill and Wanless
($89$ classes) and the $23$ P1Fs of $K_{14}$ exactly. We then give the first
\textbf{complete classification of $K_{18}$ by automorphism of order $\ge 3$}: exactly
\textbf{$531$} classes, by group order
$\{3{:}351,\,4{:}144,\,8{:}22,\,16{:}12,\,17{:}1,\,272{:}1\}$, of which only the
$\AGL(1,17)$ member was previously known; it also corrects an order-$8$ undercount
($30$, not $20$). For $K_{20}$ we find $7$ classes with
an automorphism of order $19$ and none of order $4,5,7,11,13$ or $17$. We prove a parity
theorem --- a P1F of $K_{2n}$ has an automorphism of order divisible by $4$, or a
fixed-point-free involution, only if $2n\equiv 2\pmod 4$ --- which settles every order-$2$
existence question for $K_{18}$ and $K_{20}$; only the enumeration of order-$2$ classes
remains open. As an application to \textbf{atomic Latin squares},
exactly two of the $531$ $K_{18}$ classes are atomic, giving two order-$17$ atomic squares
--- precisely the two P1Fs of $K_{18}$ of Bryant, Maenhaut and Wanless (2006) that our
census rediscovers.
\end{abstract}

\maketitle

\section{Introduction}

A \emph{1-factorisation} of $K_{2n}$ is a partition of its edge set into $2n-1$ perfect
matchings (1-factors). It is \textbf{perfect} (a P1F) if the union of every pair of
distinct 1-factors is a single Hamiltonian cycle of $K_{2n}$. P1Fs are central
objects in combinatorial design theory, with connections to atomic Latin squares,
the perfect 1-factorisation conjecture, and round-robin tournament scheduling
\cite{Wallis,MendelsohnRosa}.

Exact enumeration of P1Fs grows explosively: there is a unique P1F of $K_4,K_6,K_8$
(up to isomorphism), $5$ of $K_{12}$, $23$ of $K_{14}$, and --- the current frontier ---
$3155$ of $K_{16}$, of which $89$ have a non-trivial automorphism group, all cyclic
\cite{GillWanless}. Their computation enumerated \emph{all} $3155$ P1Fs by backtracking and
then identified automorphisms post hoc. That strategy does not extend to $K_{18}$:
the number of P1Fs is far too large to list.

Yet the P1Fs of greatest interest --- those with symmetry --- are exactly the ones that
remain accessible, because prescribing an automorphism collapses the search space
enormously. This is the classical idea behind isomorph-free generation of designs
with prescribed automorphism groups. We apply it to P1Fs:

\begin{principle}
Fix a permutation $\alpha$ of the $2n$ vertices.
Enumerate every P1F that $\alpha$ maps to itself, directly, without enumerating any
P1F that $\alpha$ does not fix.
\end{principle}

Running one representative $\alpha$ per cycle type of order $k$ and taking the union
(deduplicated up to isomorphism) yields \textbf{every} P1F of $K_{2n}$ admitting an
automorphism of order $k$. Because the cyclic-group constructions used in prior work
are a special case, this both recovers the known symmetric examples and finds the
rest.

\medskip
\noindent\textbf{Contributions.}
\begin{enumerate}
\item An orderly-generation algorithm enumerating all $\alpha$-invariant P1Fs for a
fixed $\alpha$ (\S3), with per-level isomorph rejection by the centralizer
$C(\alpha)$.
\item Exact reproduction of two known classifications, class-for-class (\S4): the
Gill--Wanless $K_{16}$ enumeration, and the $23$ P1Fs of $K_{14}$ --- the latter the
$2n\equiv 2\pmod 4$ case that validates the fixed-point-free machinery and, being
small, proves the order-$2$ classification complete.
\item New complete counts for $K_{18}$ and $K_{20}$ by automorphism order (\S5),
including $179$ classes for $K_{18}$ with an order-$4$ automorphism, and a
correction to the count obtainable by cyclic-only searches. A consequence (\S5.5):
the \emph{starter-induced} (1-rotational) P1Fs --- the entire output of the Skolem-sequence
and hill-climbing constructions of prior work --- are exactly the order-$(2n-1)$ classes,
here classified \textbf{completely} (two for $K_{18}$, seven for $K_{20}$); the genuinely new
classes are the non-1-rotational ones.
\item Two parity theorems (\S5.3): a P1F of $K_{2n}$ admits an automorphism of order
divisible by $4$ only when $2n\equiv2\pmod4$; and the fixed-point-free involution
($2^n$) requires $2n\equiv2\pmod4$. In particular $K_{16}$ and $K_{20}$ admit no such
automorphisms.
\item The order-$2$ landscape (\S5.4), now \textbf{complete for existence}: a one-fixed-factor
lemma for the two-fixed-point involution, the Galois construction settling the
two-fixed-point cases for $K_{18}$ and $K_{20}$, and an explicit $K_{18}$ P1F with a
fixed-point-free $2^9$ automorphism (Appendix~A.5) settling the last open case. Every
order-$2$ type is thus decided; only the full enumeration of order-$2$ classes is out
of reach.
\item An application to \textbf{atomic Latin squares} (\S5.6): exactly two of the $531$
P1Fs of $K_{18}$ are atomic, giving two distinct main classes of atomic Latin squares of
order $17$ --- the classical $\AGL(1,17)$ one and one from a $|\Aut|=16$, non-diagonally-cyclic
factorisation --- whereas $K_{16}$ gives none of order $15$ \cite{GillWanless}. These two are
precisely the two even-starter P1Fs of $K_{18}$ of Bryant, Maenhaut and Wanless
\cite{BryantMaenhautWanless}; our automorphism-first census independently rediscovers both,
confirmed by an explicit reconstruction-and-paratopy match. The section thus doubles as a
cross-validation against the literature.
\end{enumerate}

\section{Preliminaries}

Let $V=\{0,1,\dots,2n-1\}$ and $G=K_{2n}$. A \textbf{1-factor} (perfect matching) is a set
of $n$ pairwise-disjoint edges covering $V$; we represent it in adjacency form
$M:V\to V$ with $M(u)$ the partner of $u$. Two 1-factors $A,B$ are \textbf{perfect} if
$A\cup B$ is a single Hamiltonian cycle. A \textbf{P1F} is a set $\mathcal F$ of $2n-1$
pairwise-perfect 1-factors partitioning $E(G)$.

An \textbf{automorphism} of a P1F $\mathcal F$ is a vertex permutation $\alpha\in S_{2n}$
with $\alpha(\mathcal F)=\mathcal F$ (equivalently, $\alpha$ permutes the 1-factors of
$\mathcal F$). The automorphisms form a group $\Aut(\mathcal F)$. Two P1Fs are
\textbf{isomorphic} if a vertex permutation maps one to the other; we count isomorphism
classes.

For a permutation $\alpha$, its \textbf{cycle type} is the multiset of its cycle lengths;
its \textbf{order} is their least common multiple. The \textbf{centralizer}
$C(\alpha)=\{g\in S_{2n}: g\alpha=\alpha g\}$ acts on the set of $\alpha$-invariant
P1Fs and on partial structures; it is the symmetry group we exploit.

If $\mathcal F$ is $\alpha$-invariant then $\alpha$ permutes its $2n-1$ factors, so
$\mathcal F$ is a disjoint union of \textbf{$\alpha$-orbits of factors}
$\{M,\alpha M,\alpha^2M,\dots\}$. Each such orbit consists of pairwise-perfect,
edge-disjoint 1-factors, and the orbit lengths divide the order of $\alpha$.

\section{Method}

\subsection{Reduction to cycle types}

Conjugate permutations yield isomorphic invariant-P1F families, and conjugacy in
$S_{2n}$ is determined by cycle type. Hence to find all P1Fs admitting an
automorphism of order $k$ it suffices to:
\begin{enumerate}
\item enumerate the cycle types of order $k$ (partitions of $2n$ into parts whose lcm
is $k$);
\item for one representative $\alpha$ of each type, enumerate all $\alpha$-invariant
P1Fs; and
\item deduplicate the union up to isomorphism.
\end{enumerate}
For order $k$ the relevant cycle types are partitions of $2n$ into parts dividing
$k$ with least common multiple exactly $k$ --- typically only a few dozen, versus the
astronomically many P1Fs.

\subsection{Exact cover by factor-orbits}

Fix $\alpha$. We build an $\alpha$-invariant P1F as an \textbf{exact cover} of $E(G)$ by
$\alpha$-invariant factor-orbits. Order the edges $0,1,\dots,\binom{2n}{2}-1$
lexicographically. At each step:
\begin{enumerate}
\item let $e$ be the lowest uncovered edge;
\item generate every 1-factor $M\ni e$ that is perfect with all committed factors and
whose orbit $\{M,\alpha M,\dots\}$ is valid (members pairwise perfect and
edge-disjoint);
\item commit one such orbit and recurse; when all edges are covered, $\mathcal F$ is a
complete $\alpha$-invariant P1F.
\end{enumerate}
Candidate factors through $e$ are generated by a matching search with Hamiltonicity
pruning (a partial matching is rejected as soon as it cannot complete to a factor
perfect with a committed one).

\subsection{Isomorph rejection by orderly generation}

Two cover branches related by an element of $C(\alpha)$ produce isomorphic P1Fs, so
exploring both is wasteful. We use \textbf{orderly generation} over $C(\alpha)$:

\begin{ogrule}
At a node with committed factor-set $P$, let $G\le C(\alpha)$ be the subgroup
fixing $P$ setwise. Among the candidate factor-orbits covering the lowest uncovered
edge, keep \textbf{one representative per $G$-orbit}, and recurse into each chosen orbit
$O$ with the smaller stabilizer $\operatorname{Stab}_G(O)$.
\end{ogrule}

This is sound for counting isomorphism classes: any $g\in G$ is a bijection between
the sets of completions extending two $G$-equivalent siblings, so exploring one
representative reaches every class the others would. As the search descends, the
stabilizer chain $G\ge\operatorname{Stab}_G(O)\ge\cdots$ shrinks quickly, collapsing the
dense high-symmetry branches that defeat na\"{\i}ve pruning.

A labelling-independent canonical form (the lexicographically least serialisation
over all Hamiltonian-cycle framings of the factor set) provides the final
deduplication across cycle types and labellings, and simultaneously yields
$|\Aut(\mathcal F)|$ (as the number of framings attaining the minimum,
divided by two).

\subsection{Complexity remarks}

When $|C(\alpha)|$ is moderate the centralizer is enumerated directly and the per-level
$G$-orbit dedup is cheap. When $|C(\alpha)|$ is large --- cycle types with many fixed
points, or the fixed-point-free involution --- the centralizer is instead carried by a
strong generating set, and the per-level orbits and stabilizers are computed from it
(Schreier--Sims) without enumerating its elements, up to a size cap (of order $10^6$);
this is what completes the order-$2$ classifications of $K_{14}$ and $K_{16}$ (\S4).
Beyond the cap the dedup is disabled and the branch is walked without it --- the present
performance bottleneck (\S6). The over-cap cases are the dense fixed-point types (which,
empirically, yield no P1F, so they cost time without changing the counts) and the
two-fixed-point involution $2^{n-1}1^2$ of $K_{18}$ and $K_{20}$.

\subsection{Which orders suffice}

The method classifies P1Fs admitting an automorphism of a \emph{given} order; to obtain a
\textbf{complete} classification we must choose a set of orders whose union captures every
class of interest. The choice is not arbitrary, and rests on one elementary fact.

\begin{lemma}[sufficient orders]
A finite group contains an element of order $\ge 3$
\textbf{if and only if} it contains an element of odd prime order or an element of order
$4$.
\end{lemma}

\begin{proof}
If $g$ has order $m\ge3$ and $m$ has an odd prime divisor $q$, then $g^{m/q}$
has order $q$; if instead $m$ is a power of $2$ (hence $m\ge4$), then $g^{m/4}$ has
order $4$. Conversely an element of odd prime order or of order $4$ has order
$\ge3$.
\end{proof}

Consequently, to enumerate \textbf{every} P1F of $K_{2n}$ admitting an automorphism of
order $\ge 3$ it suffices to run the sweep over
\[
L\in\{\text{odd primes }\le 2n-1\}\cup\{4\}.
\]
For $K_{18}$ this is $L\in\{3,4,5,7,11,13,17\}$; for $K_{20}$,
$L\in\{3,4,5,7,11,13,17,19\}$. Every composite order and every higher power of $2$ is
\textbf{redundant} for completeness --- any class it could find already contains an element
of odd prime order or of order $4$, so it appears under some $L$ in the set above.
Sweeps at order $8$ (and $16,\dots$) are run only as an independent cross-check and to
expose the $2$-power \emph{element} structure, not because they add classes. This is the
precise sense in which our classification is complete ``for order $\ge 3$'': the one
order outside the set is order $2$, which is the single intractable case (\S6). It is
also why order $2$ --- and hence the full $|\Aut|>1$ classification --- cannot be
reached by adding a few more sweeps: an automorphism group that is an elementary
abelian $2$-group has no element of order $\ge 3$ at all, so it is invisible to every
$L$ in the set and surfaces only in the (intractable) order-$2$ search.

\section{Validation}

\subsection{$K_{16}$}

Running the method over the sufficient set of orders (\S3.5) --- here $\{2,3,5,7\}$ with
their composite proxies --- and deduplicating the union reproduces the \textbf{complete
Gill--Wanless catalogue} of $K_{16}$ P1Fs with a non-trivial automorphism group: all
\textbf{$89$} classes, class-for-class, with the independently computed $|\Aut|$
distribution

\begin{center}
\begin{tabular}{lccccccc}
\toprule
$|\Aut(\mathcal F)|$ & $2$ & $3$ & $5$ & $7$ & $14$ & $15$ & \textbf{total} \\
\midrule
classes & $59$ & $19$ & $5$ & $4$ & $1$ & $1$ & \textbf{$89$} \\
\bottomrule
\end{tabular}
\end{center}

\noindent matching Gill--Wanless \cite{GillWanless} exactly --- including their finding that
every such group is cyclic. The odd-prime part of this census, the $30$ classes with
$|\Aut|>2$ not a power of two, is found by orders $3,5,7$ and breaks down by
generator cycle type as:

\begin{center}
\begin{tabular}{lccc}
\toprule
generator cycle type & order & this work & Gill--Wanless \\
\midrule
$3^5 1^1$  & 3  & 19 & 19 \\
$5^3 1^1$  & 5  & 5  & 5  \\
$7^2 1^2$  & 7  & 4  & 4  \\
$14^1 1^2$ & 14 & 1  & 1  \\
$15^1 1^1$ & 15 & 1  & 1  \\
\midrule
\textbf{total} ($|\Aut|>2$) & & \textbf{30} & \textbf{30} \\
\bottomrule
\end{tabular}
\end{center}

Two features make this the strongest available check of the engine. First, it is
\textbf{complete, including order $2$}: unlike $K_{18}$ and $K_{20}$, the $K_{16}$ order-$2$
types are small enough to finish --- the fixed-point-free type $2^8$ comes back empty
directly from the enumeration, an independent computational check of the order-$2$
parity theorem (\S5.3), and the two-fixed-point type $2^7 1^2$ enumerates fully, yielding
all $59$ classes with $|\Aut|=2$ (together with the shared $|\Aut|=14$ class, which also
carries an order-$2$ element, the order-$2$ sweep returns exactly $60$ classes, with no
artificial cutoff). Second, the census carries
\textbf{no $|\Aut|$ divisible by $4$} --- the $2$-part of every automorphism group is
exactly $C_2$ --- corroborating the order-$4$ parity theorem of \S5.3 ($K_{16}$ admits no
order-$4$ automorphism) and showing empirically that $K_{16}$ has no Klein-$V_4$-symmetric
P1F either.

The independently computed automorphism group orders match in every class. The dense
type $3^5 1^1$ (centralizer order $29160$), on which a na\"{\i}ve symmetry-breaking search
stalls, completes in seconds under orderly generation. This agreement with the
published authority --- not merely on the total, but on the full distribution --- gives
strong confidence in the implementation.

\subsection{$K_{14}$ --- the fixed-point-free case}

$K_{14}$ has exactly $23$ P1Fs. Two of them are asymmetric ($|\Aut|=1$) and
hence invisible to any automorphism-first method; the remaining $21$ have a
non-trivial automorphism group, with the distribution
$\{2{:}3,\,3{:}5,\,4{:}1,\,6{:}5,\,12{:}5,\,84{:}1,\,156{:}1\}$. Our method reproduces
this distribution \textbf{exactly}, cross-checked against an independent complete
enumeration:

\begin{center}
\begin{tabular}{lcl}
\toprule
automorphism order & classes & $|\Aut|$ distribution \\
\midrule
2 & 16 & $\{2{:}3,\,4{:}1,\,6{:}5,\,12{:}5,\,84{:}1,\,156{:}1\}$ \\
3 & 17 & $\{3{:}5,\,6{:}5,\,12{:}5,\,84{:}1,\,156{:}1\}$ \\
7 & 1  & $\{84{:}1\}$ \\
13 & 1 & $\{156{:}1\}$ \\
\midrule
\textbf{union} ($|\Aut|>1$) & \textbf{21} & $\{2{:}3,\,3{:}5,\,4{:}1,\,6{:}5,\,12{:}5,\,84{:}1,\,156{:}1\}$ \\
\bottomrule
\end{tabular}
\end{center}

This validation is qualitatively stronger than the $K_{16}$ one in two respects.

\medskip
\noindent\emph{It exercises the structure of $K_{18}$.} Because $14\equiv 2\pmod 4$, like $18$ and
unlike $16$, the order-$2$ and order-$4$ cycle types of $K_{14}$ include
\textbf{fixed-point-free} ones ($2^7$ and $4^3 2^1$) --- precisely the types responsible for
the $K_{18}$ results in \S5 and structurally absent from $K_{16}$. The identical engine
producing the exact known $K_{14}$ counts is therefore direct evidence for the
correctness of the $K_{18}$ machinery, not merely of the $N$-generic core.

\medskip
\noindent\emph{It closes the order-$2$ case.} For $K_{14}$ the fixed-point-free involution type
$2^7$ has centralizer of order $2^7\cdot 7! = 645{,}120$, small enough to enumerate,
so order $2$ \textbf{completes} --- yielding the proven count of $16$ in a few seconds. This
is the case that is intractable for the larger sizes (where $2^8$, $2^9$, $\dots$ have
centralizers too large to enumerate; cf.\ \S6), and $K_{14}$ confirms that the
limitation there is one of scale, not of correctness.

\medskip
\noindent\emph{A remark on composite orders.} By Cauchy's theorem a group of order divisible by a
prime $p$ contains an element of order $p$, so for \textbf{prime} orders the count of
classes our method finds equals the number whose $|\Aut|$ is divisible by that
prime --- making orders $2,3,7,13$ above clean checks. This fails for \textbf{composite}
orders: a group whose order is divisible by $4$ need not contain an element of order
$4$ (its Sylow $2$-subgroup may be elementary abelian, or the group may be $A_4$).
Indeed $K_{14}$ has $8$ classes with $4\mid|\Aut|$ but only $5$ admitting an
order-$4$ \emph{automorphism} ($\{12{:}4,\,156{:}1\}$); the three excluded are the
$|\Aut|=4$ (Klein), one $|\Aut|=12$ ($A_4$, found instead by order
$3$), and the $|\Aut|=84$ class. The order-$4$ counts reported throughout this
paper are counts of P1Fs \emph{admitting an order-$4$ automorphism}, i.e.\ with an order-$4$
group element.

\section{New results}

\subsection{$K_{18}$}

$K_{18}$ admits no complete published enumeration. The only P1F of $K_{18}$ on record
is the algebraic construction for orders $q+1$ ($q$ a prime power), here $q=17$,
whose automorphism group is the affine group $\AGL(1,17)$ of order
$17\cdot 16 = 272$. Our method recovers it. Running the sufficient set of orders
$L\in\{3,4,5,7,11,13,17\}$ (\S3.5) and deduplicating the union yields the \textbf{complete
classification of $K_{18}$ by automorphism of order $\ge 3$}: exactly \textbf{$531$}
isomorphism classes, with the distribution

\begin{center}
\begin{tabular}{lccccccc}
\toprule
$|\Aut(\mathcal F)|$ & $3$ & $4$ & $8$ & $16$ & $17$ & $272$ & \textbf{total} \\
\midrule
classes & $351$ & $144$ & $22$ & $12$ & $1$ & $1$ & \textbf{$531$} \\
\bottomrule
\end{tabular}
\end{center}

To our knowledge this is the first complete such classification for $K_{18}$; every
count in it except the single $\AGL(1,17)$ member is new. The orders
$5,7,11,13$ contribute nothing --- no $|\Aut|$ above is divisible by any of them,
which by Cauchy's theorem is exactly the statement that those sweeps are empty --- so the
symmetry present in $K_{18}$ P1Fs is confined to the primes $3,17$ and the powers of
$2$.\footnote{The census is the union of the per-order sweeps and is best read by
$|\Aut|$ value, but it is worth recording the by-order contributions, since one
of them corrects a prior undercount. The order-$4$ sweep finds the $179$ classes whose
group contains an order-$4$ element ($\{4{:}144,\,8{:}22,\,16{:}12,\,272{:}1\}$ --- all
$22$ groups of order $8$ and all $12$ of order $16$ here contain such an element); the
order-$8$ sweep finds the $30$ whose group contains an order-$8$ element
($\{8{:}17,\,16{:}12,\,272{:}1\}$); and the order-$17$ sweep finds the $2$ with
$17\mid|\Aut|$. The order-$8$ figures $22$ (groups of order $8$) versus $17$
(those with an order-$8$ element) differ because only the cyclic $C_8$ groups contain
an order-$8$ element --- the other $5$ groups of order $8$ (e.g.\ $C_4\times C_2$, $D_4$)
have order-$4$ elements only, so they are found by the order-$4$ sweep, not order-$8$.
This also underlies a correction: a search that builds the automorphism cyclically
(fixing a vertex) structurally misses P1Fs whose order-$8$ automorphism is
fixed-point-free (cycle type $8^2 2^1$) and reports only $20$ classes for order $8$;
the automorphism-first method, placing one representative per cycle type including the
fixed-point-free ones, gives the correct $30$.} The $144$ classes at $|\Aut|=4$ are all cyclic ($C_4$); the
$|\Aut|=272$ member is the $\AGL(1,17)$ factorisation (Appendix~A.1),
the $|\Aut|=17$ member the cyclic Skolem-starter one (Appendix~A.4, \S5.5).

Exactly one of the $531$ classes --- one with $|\Aut|=16$ --- additionally carries
a fixed-point-free involution of type $2^9$; this is the witness that settles the last
open order-$2$ case (\S5.4, Appendix~A.5).

\subsection{$K_{20}$}

$K_{20}=K_{19+1}$ likewise has a single recorded P1F, from $\AGL(1,19)$ of
order $19\cdot 18=342$; our method recovers it. Complete counts so far:

\begin{center}
\begin{tabular}{lcl}
\toprule
automorphism order & classes & $|\Aut|$ distribution \\
\midrule
19 & 7 & $\{19{:}3,\ 57{:}1,\ 171{:}2,\ 342{:}1\}$ \\
17 & 0 & --- \\
13 & 0 & --- \\
11 & 0 & --- \\
7  & 0 & --- \\
5  & 0 & --- \\
4  & 0 & --- \\
\bottomrule
\end{tabular}
\end{center}

No P1F of $K_{20}$ admits an automorphism of order $4,5,7,11,13,$ or $17$. The seven
order-$19$ classes all lie in the $\AGL(1,19)$ family (orders dividing
$342=2\cdot 3^2\cdot 19$). The count for order $3$ is in progress (\S6).

\subsection{A parity observation}

The fixed-point-free order-$4$ cycle types ($4^4$ for $K_{16}$, $4^5$ and $4^4 2^2$
for $K_{20}$, etc.) admit no valid factor-orbit through the first edge, so they
contribute nothing; for the sizes with order-$4$ P1Fs the contributing types are
those with two fixed points ($4^3 1^2$ for $K_{14}$, $4^4 1^2$ for $K_{18}$). Across
the sizes accessible to us this gives
\[
\#\{K_{14}\}=5,\quad \#\{K_{16}\}=0,\quad \#\{K_{18}\}=179,\quad \#\{K_{20}\}=0,
\]
where $\#\{K_{2n}\}$ denotes the number of P1Fs of $K_{2n}$ admitting an order-$4$
automorphism. The two sizes with such P1Fs are exactly the two with
$2n\equiv 2\pmod 4$ ($K_{14},K_{18}$); the two without are exactly those with
$2n\equiv 0$ ($K_{16},K_{20}$). This is now a theorem:

\begin{theorem}[order-$4k$ parity]
A P1F of $K_{2n}$ ($2n\ge8$) admits an
automorphism whose order is \textbf{divisible by $4$} (orders $4,8,12,\dots$) only if
$2n\equiv 2\pmod 4$ (equivalently $n$ odd). In particular $K_{16}$ and $K_{20}$ admit
no P1F automorphism of any order divisible by $4$.
\end{theorem}

\begin{proof}
Let $\beta$ have order $4s$ and put $\gamma=\beta^{2s}$ (an involution). A
\textbf{lemma} bounds the fixed points of any involution automorphism: $\gamma$ acts on any
$\gamma$-invariant factor-pair's Hamiltonian cycle as an involution of $D_{2n}$, which
fixes $0$ or $2$ vertices --- and such a pair always exists (the number of $\gamma$-fixed
factors is odd, hence $\ge1$; pair two fixed factors, or a fixed factor's swapped
partners). Now $\gamma$'s moved points lie in $\beta$-cycles of length divisible by $4$
(those with $L\mid4s,\,L\nmid2s$), so the moved count is $\equiv0\pmod4$ and the fixed
count $f\equiv2n\pmod4$; with $f\le2$ even, $f\in\{0,2\}$. If $f=2$ then $2n\equiv2$, i.e.\
$n$ odd. If $f=0$ then $\gamma$ is the fixed-point-free $2^n$ with $2n\equiv0$ ($n$ even),
which the order-2 theorem below forbids. Hence $n$ is odd.
\end{proof}

Full details and computational validation (orders $4,8,12$) are in the supplementary
material. Orders $\equiv2\pmod4$ do not reduce this way (they meet the open
mixed type $2^{n-1}1^2$). The \textbf{order-$2$ theorem} is the engine of the $f=0$ step:

\begin{theorem}[order-2 parity]
For $2n\ge6$, $K_{2n}$ has a P1F invariant under a
fixed-point-free involution (cycle type $2^n$) only if $2n\equiv2\pmod4$. In
particular $K_{16}$ has no $2^8$-invariant and $K_{20}$ no $2^{10}$-invariant P1F.
\end{theorem}

\begin{proof}[Proof sketch]
The $2n-1$ factors split under $\alpha$ into fixed factors and swapped
pairs $\{G,\alpha G\}$. For a swapped pair, $G\cup\alpha G$ is a Hamiltonian $2n$-cycle
$H$ on which $\alpha$ acts as a fixed-point-free automorphism; as $G,\alpha G$ carry no
$\alpha$-fixed (diagonal) edge, that automorphism fixes no edge of $H$ and is therefore
the antipodal rotation $\rho$. The antipodal rotation of a $2n$-cycle swaps its two
$1$-factors iff $n$ is odd; since $\alpha$ swaps $G$ and $\alpha G$, $n$ must be odd.
Hence for $n$ even every factor is $\alpha$-fixed, and the $n$ diagonals (each in a fixed
factor, each fixed factor holding $\equiv n\equiv0\pmod2$ of them, and any two fixed
factors' union holding $0$ or $2$) cannot sum to $n\ge4$ --- contradiction.
\end{proof}

Both parity statements are now theorems, and the order-4 one generalises to \textbf{every
order divisible by $4$} (same proof with $\beta^2$ replaced by $\beta^{\,\mathrm{ord}/2}$):
no P1F of $K_{16}$ or $K_{20}$ admits an automorphism of order $4,8,12,16,\dots$.

\subsection{The order-2 landscape}

The theorems of \S5.3 settle the \textbf{fixed-point-free} involution type $2^n$ for
$2n\equiv0\pmod4$ ($K_{16}\,2^8$ and $K_{20}\,2^{10}$ are empty). The other involution
type --- by the $\le 2$-fixed-points lemma these are the only two --- is the
\textbf{two-fixed-point} type $2^{n-1}1^2$. We settle its existence and pin its structure.

\begin{lemma}[one fixed factor]
If a P1F of $K_{2n}$ has an automorphism $\alpha$ of cycle
type $2^{n-1}1^2$ (fixing exactly two vertices $u,v$), then $\alpha$ fixes \textbf{exactly one}
factor: the \emph{diagonal factor} $D$ consisting of the edge $\{u,v\}$ together with the
$n-1$ transposition pairs (these $n$ $\alpha$-fixed edges form one perfect matching).
\end{lemma}

\begin{proof}
If two distinct factors $F,G$ were both $\alpha$-fixed, then $H=F\cup G$ is a
Hamiltonian cycle with $\alpha(H)=H$; as $H$ spans $V$, $\alpha|_H$ is the cycle
involution fixing exactly $u,v$ --- the reflection through them. On an even cycle the
unique proper $2$-edge-colouring gives the two edges at each fixed vertex opposite
colours, and that reflection swaps them, hence swaps $F$ and $G$ --- contradicting both
being fixed. So at most one fixed factor; and an involution acting on the odd number
$2n-1$ of factors fixes an odd number, so exactly one. It carries all $n$ $\alpha$-fixed
edges, i.e.\ equals $D$.
\end{proof}

The lemma was validated against direct enumeration: for $K_6,K_8,K_{10},K_{12}$ the
number of $2^{n-1}1^2$-invariant P1Fs is $2,16,96,1536$ respectively, and in every one
the fixed-factor count is exactly $1$.

\medskip
\noindent\textbf{Existence (Galois construction).} The two-fixed-point case is non-empty for both our
sizes. The Galois P1F $GK_{p+1}$ over $\GF(p)$ ($p$ prime; vertices
$\GF(p)\cup\{\infty\}$, factor $F_k=\{\infty,k\}\cup\{\{k+i,k-i\}:1\le i\le(p-1)/2\}$)
admits the involution $x\mapsto -x$, fixing $0$ and $\infty$, of cycle type
$2^{(p-1)/2}1^2$. For $p=17$ this is the $\AGL(1,17)$ factorisation of $K_{18}$
(Appendix~A.1), type $2^8 1^2$; for $p=19$ the $\AGL(1,19)$ factorisation of
$K_{20}$ (Appendix~A.3), type $2^9 1^2$. Hence \textbf{$K_{18}$ and $K_{20}$ both admit a P1F
with a two-fixed-point involution.} The existence question is thus settled by
construction; only the \emph{complete enumeration} of these classes is out of reach, their
centralizers ($2^{n-1}(n-1)!\cdot2$) being far past the enumeration cap (\S6).

\medskip
\noindent\textbf{The last fixed-point-free case, resolved.} The one remaining order-$2$ existence
question --- the fixed-point-free type $2^9$ of $K_{18}$ --- we settle \textbf{affirmatively}:
$K_{18}$ admits a perfect $1$-factorisation invariant under a fixed-point-free involution
of type $2^9$. This is notable because the type resists direct attack on all three of the
usual fronts: the Galois construction yields only the two-fixed-point type, never
fixed-point-free; the parity theorem of \S5.3 does not apply ($18\equiv2\pmod4$, and
$K_{14}\,2^7$ being non-empty shows there is no parity obstruction to exploit); and a
direct exhaustive search over the type is infeasible (its centralizer
$C(\alpha)=C_2\wr S_9$ has order $2^9\cdot 9!\approx 1.9\times10^8$, far past the
enumeration cap, and the search tree is estimated at $\sim 10^{16}$ nodes).

The witness is obtained not by searching the $2^9$ type directly but as a by-product of
the tractable \textbf{order-$4$} enumeration (\S5.1). Any P1F whose automorphism group has order
divisible by $4$ and contains a $2^9$ element is found by the order-$4$ sweep; exactly one
of the $179$ order-$4$ classes --- a class $\mathcal F^\ast$ with $|\Aut|=16$ --- is of
this kind. Its group contains the fixed-point-free involution
\[
\alpha=(0\,1)(2\,4)(3\,5)(6\,8)(7\,9)(10\,12)(11\,13)(14\,16)(15\,17),
\]
which permutes the $17$ factors of $\mathcal F^\ast$ as $9$ fixed factors and $4$ swapped
pairs --- exactly the structure forced below. The factorisation is listed in Appendix~A.5,
was verified independently (all $\binom{17}{2}=136$ factor-pairs are Hamiltonian; $\alpha$
is a genuine automorphism), and was reproduced as a byte-identical canonical matrix by a
separate run. Its $9$ fixed factors realise the $K_{10}$-lift described next.

This realises the one structural handle on the type: the $n=9$ fixed factors of any
$2^9$-invariant P1F of $K_{18}$ form a perfect $1$-factorisation of $K_{10}$, and
\textbf{$K_{10}$ has a unique P1F up to isomorphism} (verified: $864$ labellings with the first
factor fixed, one class). The example $\mathcal F^\ast$ exhibits a lift of that single
object to a $2^9$-invariant P1F of $K_{18}$, so the lift exists.

\medskip
\noindent\textbf{The order-$2$ existence landscape is now complete.} Combining \S5.3 and the above:

\begin{center}
\begin{tabular}{llll}
\toprule
size & type & cycle type & status \\
\midrule
$K_{16}$ & fixed-point-free & $2^8$    & \textbf{empty} (\S5.3) \\
$K_{20}$ & fixed-point-free & $2^{10}$ & \textbf{empty} (\S5.3) \\
$K_{18}$ & fixed-point-free & $2^9$    & \textbf{non-empty} (Appendix~A.5) \\
$K_{18}$ & two-fixed-point  & $2^8 1^2$ & \textbf{non-empty} (Galois, A.1) \\
$K_{20}$ & two-fixed-point  & $2^9 1^2$ & \textbf{non-empty} (Galois, A.3) \\
\bottomrule
\end{tabular}
\end{center}

Every order-$2$ involution type for $K_{18}$ and $K_{20}$ is decided. What remains beyond
the present method is only the \emph{complete enumeration} of the order-$2$ classes --- equivalently
the P1Fs whose automorphism group is an elementary-abelian $2$-group, invisible to any
order-$\ge3$ (or order-$4$) sweep --- whose centralizers exceed the enumeration cap (\S6).

\subsection{The 1-rotational (starter-induced) subfamily and prior constructions}

Every previously recorded P1F of $K_{18}$ and $K_{20}$ arises from a \emph{starter}, and it is
worth pinning down precisely where these sit in our classification. A P1F of $K_{2n}$ on the
vertex set $\mathbb{Z}_{2n-1}\cup\{\infty\}$ is \textbf{1-rotational} if the rotation
$\rho:x\mapsto x+1\pmod{2n-1}$ (with $\rho(\infty)=\infty$) is an automorphism; equivalently
it is the cyclic development of a \emph{starter} --- a set of $n-1$ pairs partitioning
$\mathbb{Z}_{2n-1}\setminus\{0\}$ whose differences cover $\mathbb{Z}_{2n-1}\setminus\{0\}$,
together with the pair $\{0,\infty\}$. The rotation $\rho$ has order $2n-1$ and cycle type
$(2n-1)^1\,1^1$.

For both our sizes $2n-1\in\{17,19\}$ is \textbf{prime}, which makes the correspondence exact:

\begin{observation}
For $K_{2n}$ with $2n-1$ prime, a P1F is starter-induced (1-rotational)
\textbf{iff} $(2n-1)\mid|\Aut|$.
\end{observation}

Indeed, any automorphism of prime order $2n-1$ permutes the $2n=(2n-1)+1$ vertices as a single
$(2n-1)$-cycle and one fixed point (no other cycle type of that order fits $2n$ points), so ---
after relabelling --- it is the rotation $\rho$ and the P1F is 1-rotational; conversely $\rho$
contributes the factor $2n-1$ to $|\Aut|$. Hence the \emph{entire} starter-induced subfamily
is exactly the \textbf{order-$(2n-1)$} column of our tables, which our method classifies
\textbf{completely}:
\begin{itemize}
\item \textbf{$K_{18}$:} exactly \textbf{two} 1-rotational P1Fs (\S5.1), $|\Aut|\in\{17,\,272\}$;
\item \textbf{$K_{20}$:} exactly \textbf{seven} 1-rotational P1Fs (\S5.2),
$|\Aut|\in\{19,\,57,\,171,\,342\}$.
\end{itemize}

These complete counts subsume the prior literature. The Galois factorisations $GK_{p+1}$ are the
maximal-symmetry members ($|\Aut|=\AGL(1,p)$, i.e.\ $272$ for $K_{18}$ and $342$
for $K_{20}$; Appendices~A.1, A.3). Starter constructions from \textbf{Skolem sequences}
\cite{PikeShalaby,Pike} produce further 1-rotational members --- for example the order-$8$ Skolem
sequence $(1,1,3,7,8,3,2,6,2,5,7,4,8,6,5,4)$ develops into the $|\Aut|=17$ P1F of
$K_{18}$ (Appendix~A.4) --- but, being constructions, they exhibit individual examples rather than
a classification. \textbf{Hill-climbing} searches \cite{DinitzStinson} likewise generate starter-induced
P1Fs but cannot certify completeness, and are structurally biased \emph{against} the symmetric classes
that are our subject: a randomised generator reaches an isomorphism class with probability
roughly proportional to $1/|\Aut|$, so the high-symmetry P1Fs we enumerate are precisely
the ones such methods almost never find.

\medskip
\noindent\textbf{Consequence for the contribution.} The order-$(2n-1)$ rows therefore \emph{recover} known
territory --- the Galois and Skolem-starter factorisations are not new --- but ours is the first
enumeration to prove them \textbf{complete} (two for $K_{18}$, seven for $K_{20}$, with full
$|\Aut|$ distributions). The genuinely new P1Fs are the \textbf{non-1-rotational} ones: every
class with $|\Aut|\ge3$ whose order is coprime to $2n-1$ --- all $179$ order-$4$ classes of
$K_{18}$ among them --- has \emph{no} starter and lies outside the reach of any cyclic, Skolem-starter,
or hill-climbing method.

\subsection{Atomic Latin squares}

A prominent motivation for perfect 1-factorisations of complete graphs is their link to
\emph{atomic} Latin squares --- the ``indivisible'' squares whose structure mimics that of
the cyclic group of prime order, and of which no example of non-prime-power order is known.
We determine which of our $K_{18}$ factorisations produce atomic squares.

Following Gill--Wanless \cite{GillWanless}, from a 1-factorisation $\mathcal F$ of $K_{2n}$
and a vertex $j$ one builds a Latin square $I(\mathcal F,j)$ of order $2n-1$: start from the
symmetric \emph{unipotent} square $U(\mathcal F)$ of order $2n$ with $U[i][i']=k$ when edge
$ii'$ lies in the $k$-th factor; replace each diagonal cell $(i,i)$ by $U[i][j]$; and delete
row and column $j$. The result is symmetric and idempotent. A Latin square is
\textbf{row-Hamiltonian} (respectively \textbf{symbol-}, \textbf{column-Hamiltonian}) if
every cycle formed by a pair of its rows (respectively symbols, columns) has full length,
and \textbf{atomic} if all three hold --- equivalently, if all six conjugates are
row-Hamiltonian. Two facts frame the computation: $\mathcal F$ is perfect \textbf{iff} every
$I(\mathcal F,j)$ is symbol-Hamiltonian; and, since $I(\mathcal F,j)$ is symmetric (so column-
and row-Hamiltonicity coincide), for a P1F \textbf{$I(\mathcal F,j)$ is atomic iff it is
row-Hamiltonian}.

For $K_{16}$ this yields nothing --- Gill--Wanless proved that no symmetric atomic Latin
square of order $15$ exists. But for $K_{18}$ the order $2n-1=17$ is \textbf{prime}, and
atomic squares of prime order can exist, so the question is genuine. Testing all $531$
classes at every vertex $j$ gives:

\begin{quote}
Exactly \textbf{two} of the $531$ P1Fs of $K_{18}$ yield an atomic Latin square of order
$17$, and the two squares lie in \textbf{distinct main classes}: the $\AGL(1,17)$
factorisation ($|\Aut|=272$), giving the classical atomic square; and one class with
$|\Aut|=16$, giving a second, distinct atomic square (Appendix~A.6).
\end{quote}

The computation is self-validating through the perfect-iff-symbol-Hamiltonian equivalence:
all $531\times 18$ squares $I(\mathcal F,j)$ are symbol-Hamiltonian, as they must be.
Distinctness of the two atomic squares was verified directly --- no conjugate of one is
isotopic to the other.

These two classes are exactly the two P1Fs of $K_{18}$ recorded by Bryant, Maenhaut and
Wanless \cite{BryantMaenhautWanless}, who show that for every prime $p\ge 11$ there are two
even-starter-induced P1Fs of $K_{p+1}$: the well-known patterned factorisation $GK_{p+1}$,
and one further class. For $p=17$ the patterned factorisation is our $\AGL(1,17)$ member
($|\Aut|=272$), while the second even starter $E_B$ produces our $|\Aut|=16$ class. We
confirmed this correspondence explicitly: reconstructing both even-starter factorisations of
$K_{18}$ from \cite{BryantMaenhautWanless} and folding them to order-$17$ squares yields
squares \textbf{paratopic} to our two (autoparatopism-group orders $\nu=27744$ and $\nu=96$
respectively; \S A.6). Our automorphism-first census therefore \textbf{independently
rediscovers} their two factorisations, and this section is best read as a cross-validation of
the classification against the literature rather than a new-square claim. The $|\Aut|=16$
class is genuinely not diagonally cyclic (since $17\nmid16$, it admits no automorphism of
order $17$), placing it --- as \cite{BryantMaenhautWanless} observe --- outside the
cyclotomic-orthomorphism and diagonally-cyclic constructions; note that $2$ is not a
primitive root modulo $17$, so the prime-order atomic family of \cite{BryantMaenhautWanless}
(which requires $2$ primitive) does not cover order $17$. For $K_{20}$ ($2n-1=19$, again
prime) the same analysis is pending completion of the order-$\ge 3$ sweep (\S6).

\section{Status of results, reproducibility, and open problems}

\noindent\textbf{Status.} The $K_{16}$ and $K_{14}$ validations (the latter including a \emph{complete}
order-$2$ classification) are complete and reproducible. For $K_{18}$ the full
sufficient sweep $L\in\{3,4,5,7,11,13,17\}$ (\S3.5) has been run and deduplicated,
giving the \textbf{complete classification by automorphism of order $\ge 3$} --- the $531$
classes of \S5.1 --- so this result is final, not partial. For $K_{20}$ the counts for
orders $4,5,7,11,13,17,19$ are complete and reproducible; the $K_{20}$ count for order
$3$ (a dense over-cap type, below) is the one order-$\ge3$ computation still in
progress and will appear in a later revision.

\medskip
\noindent\textbf{The order-2 / dense-type bottleneck.} Two situations currently lack per-level
isomorph rejection: order-$2$ (the fixed-point-free type $2^n$) and cycle types with
many fixed points (e.g.\ $K_{20}$, order $3$, type $3^5 1^5$). Both have a centralizer
too large to enumerate. That this is a limitation of \emph{scale} rather than of
correctness is shown by $K_{14}$ (\S4.2): there the fixed-point-free involution type
$2^7$ has an enumerable centralizer, order $2$ completes, and the method returns the
exact known count of $16$. The obstruction at the larger sizes is only that $2^8$,
$2^9$, $\dots$ have centralizers beyond the enumeration cap. Completing the order-$2$
enumeration for $K_{18}$ and $K_{20}$ --- equivalently, accounting for the classes whose
automorphism group is a $2$-group --- is beyond the present method and remains open.
The two \emph{fully} fixed-point-free over-cap types are now settled by the order-2 parity
theorem (\S5.3): $K_{16}\,2^8$ and $K_{20}\,2^{10}$ are \textbf{empty}. The two-fixed-point
over-cap types ($K_{18}\,2^8 1^2$, $K_{20}\,2^9 1^2$) are settled for \emph{existence} by the
Galois construction (\S5.4), and the last fixed-point-free type, $K_{18}\,2^9$, is settled
\textbf{non-empty} by the explicit example of Appendix~A.5 (recovered within the order-$4$
sweep). \textbf{Every order-$2$ existence question for $K_{18}$ and $K_{20}$ is therefore now
decided} (\S5.4); what remains beyond the method is only the \emph{complete enumeration} of
these order-$2$ classes, whose centralizers exceed the enumeration cap.

\medskip
\noindent\textbf{Reproducibility.} The implementation is a single self-contained module per order
$N=2n$ (changing only $N$); the search, canonical form, and orbit machinery are
$N$-generic. [TODO: repository URL, commit hash, machine specifications, and per-run
wall-clock times.]

\medskip
\noindent\textbf{Open problems.}
\begin{enumerate}
\item \emph{[Resolved (\S5.3): proved for both order $2$ and order $4$.]} Prove (or refute) the $2n\equiv 2\pmod 4$ observation for order-$4$ automorphisms.
\item Complete the $K_{18}$ and $K_{20}$ classifications. What is complete is the
classification by automorphism of \textbf{order $\ge 3$} (for both sizes); equivalently,
every P1F whose automorphism group is not an elementary-abelian $2$-group is
accounted for. This is deliberately weaker than ``$|\Aut|\ge3$'': a P1F whose
automorphism group is an elementary-abelian $2$-group of order $\ge4$ --- for instance
the Klein four-group $V_4$, or $C_2^3$ --- has $|\Aut|\ge3$ yet contains \emph{no}
element of order $\ge3$, so it is invisible to every order-$\ge3$ sweep and would
surface only in the intractable order-$2$ search; whether any such $K_{18}/K_{20}$
classes exist is itself open. By contrast any automorphism group whose order has an odd
prime factor (by Cauchy's theorem) or that contains an element of order $4$ or $8$ \emph{is}
found --- so, e.g., the $144$ classes we report at $|\Aut|=4$ are all cyclic
($C_4$). The residue is the order-$2$ part, now settled for \emph{existence} across
the board: the fixed-point-free types are $K_{16}\,2^8$, $K_{20}\,2^{10}$ empty (\S5.3)
and $K_{18}\,2^9$ non-empty (Appendix~A.5); the two-fixed-point types are non-empty
(Galois, \S5.4). What remains open is only the complete \emph{enumeration} of these order-$2$
classes --- and hence the full $|\Aut|>1$ classification of $K_{18}$ and $K_{20}$ ---
which is out of reach of the present method.
\item Atomic Latin squares. For $K_{18}$ this is now answered (\S5.6): exactly two classes
are atomic, giving two distinct main classes of atomic squares of order $17$, and both are
identified with the two even-starter P1Fs of $K_{18}$ of Bryant, Maenhaut and Wanless
\cite{BryantMaenhautWanless} (confirmed by reconstruction and paratopy). What remains is to
run the analogous $K_{20}$ computation ($2n-1=19$ prime) once its order-$\ge 3$ sweep
completes, and to see whether it likewise recovers known order-$19$ atomic squares or
produces anything outside the current families.
\end{enumerate}

\section*{Acknowledgements}

Software implementation, computational experiments, and drafting assistance were
carried out with the help of Anthropic's Claude (Claude Code). The authors directed
the research, verified the results, and are solely responsible for the content. We
thank Ian M. Wanless for making the $K_{16}$ catalogue and constructions publicly
available, which served as our validation oracle.

\begin{thebibliography}{99}

\bibitem{GillWanless} M. J. Gill and I. M. Wanless. \emph{Perfect 1-factorisations of $K_{16}$.} Bulletin of
the Australian Mathematical Society \textbf{101} (2020). arXiv:1905.07535.

\bibitem{Wallis} W. D. Wallis. \emph{One-Factorizations.} Mathematics and its Applications, Kluwer, 1997.

\bibitem{MendelsohnRosa} E. Mendelsohn and A. Rosa. \emph{One-factorizations of the complete graph --- a survey.}
Journal of Graph Theory \textbf{9} (1985), 43--65.

\bibitem{McKay} B. D. McKay. \emph{Isomorph-free exhaustive generation.} Journal of Algorithms \textbf{26}
(1998), 306--324.

\bibitem{Read} R. C. Read. \emph{Every one a winner, or how to avoid isomorphism search when
cataloguing combinatorial configurations.} Annals of Discrete Mathematics \textbf{2}
(1978), 107--120.

\bibitem{DGM} J. H. Dinitz, D. K. Garnick, and B. D. McKay. \emph{There are 526,915,620 nonisomorphic
one-factorizations of $K_{12}$.} Journal of Combinatorial Designs \textbf{2} (1994).

\bibitem{MeszkaRosa} M. Meszka and A. Rosa. \emph{Perfect 1-factorisations.} (and related work on $k$-cycle
free one-factorisations).

\bibitem{PikeShalaby} D. A. Pike and N. Shalaby. \emph{The use of Skolem sequences to generate perfect
one-factorizations.} Ars Combinatoria \textbf{59} (2001), 153--159.

\bibitem{Pike} D. A. Pike. \emph{Non-isomorphic perfect one-factorizations from Skolem sequences.} Journal
of Combinatorial Mathematics and Combinatorial Computing \textbf{44} (2003). [page range to
verify]

\bibitem{DinitzStinson} J. H. Dinitz and D. R. Stinson. \emph{Hill-climbing algorithms for the construction of
combinatorial designs.} Annals of Discrete Mathematics \textbf{26} (1985), 101--125.
[details to verify]

\bibitem{MaenhautWanless} B. M. Maenhaut and I. M. Wanless. \emph{Atomic Latin squares of order eleven.}
Journal of Combinatorial Designs \textbf{12} (2004), 12--34.

\bibitem{WanlessCyclotomic} I. M. Wanless. \emph{Atomic Latin squares based on cyclotomic orthomorphisms.}
Electronic Journal of Combinatorics \textbf{12} (2005), \#R22.

\bibitem{BryantMaenhautWanless} D. Bryant, B. M. Maenhaut and I. M. Wanless. \emph{New families of
atomic Latin squares and perfect one-factorisations.} Journal of Combinatorial Theory, Series A
\textbf{113} (2006), 608--624.

\end{thebibliography}

\appendix

\section{Key result examples}

Each example below lists a perfect 1-factorisation as its $2n-1$ 1-factors
$F_1,\dots,F_{2n-1}$ on the vertex set $\{0,1,\dots,2n-1\}$; a 1-factor is written as
its $n$ edges $(a,b)$. By definition the union of any two factors $F_i\cup F_j$ is a
single Hamiltonian cycle. The automorphism order $|\Aut|$ is as computed and
independently verified. (These are extracted verbatim from the solver output and are
reproducible from the repository result files.)

\subsection{$K_{18}$, $|\Aut|=272$ --- the $\AGL(1,17)$ factorisation}

This is the algebraic P1F of $K_{18}=K_{17+1}$; its automorphism group is the affine
group $x\mapsto ax+b$ over $\GF(17)$ (vertex $17$ playing $\infty$), of order
$17\cdot16=272$. Our method recovers it among the order-4 classes (since $4\mid 272$).

{\footnotesize\begin{verbatim}
F1  (0,1)(2,3)(4,5)(6,7)(8,9)(10,11)(12,13)(14,15)(16,17)
F2  (0,2)(1,4)(3,6)(5,8)(7,10)(9,12)(11,14)(13,16)(15,17)
F3  (0,3)(1,5)(2,7)(4,9)(6,11)(8,13)(10,15)(12,17)(14,16)
F4  (0,4)(1,8)(2,6)(3,10)(5,12)(7,14)(9,16)(11,17)(13,15)
F5  (0,5)(1,9)(2,11)(3,7)(4,13)(6,15)(8,17)(10,16)(12,14)
F6  (0,6)(1,12)(2,10)(3,14)(4,8)(5,16)(7,17)(9,15)(11,13)
F7  (0,7)(1,13)(2,15)(3,11)(4,17)(5,9)(6,16)(8,14)(10,12)
F8  (0,8)(1,16)(2,14)(3,17)(4,12)(5,15)(6,10)(7,13)(9,11)
F9  (0,9)(1,17)(2,16)(3,15)(4,14)(5,13)(6,12)(7,11)(8,10)
F10 (0,10)(1,15)(2,17)(3,13)(4,16)(5,11)(6,14)(7,9)(8,12)
F11 (0,11)(1,14)(2,12)(3,16)(4,10)(5,17)(6,8)(7,15)(9,13)
F12 (0,12)(1,11)(2,13)(3,9)(4,15)(5,7)(6,17)(8,16)(10,14)
F13 (0,13)(1,10)(2,8)(3,12)(4,6)(5,14)(7,16)(9,17)(11,15)
F14 (0,14)(1,7)(2,9)(3,5)(4,11)(6,13)(8,15)(10,17)(12,16)
F15 (0,15)(1,6)(2,4)(3,8)(5,10)(7,12)(9,14)(11,16)(13,17)
F16 (0,16)(1,3)(2,5)(4,7)(6,9)(8,11)(10,13)(12,15)(14,17)
F17 (0,17)(1,2)(3,4)(5,6)(7,8)(9,10)(11,12)(13,14)(15,16)
\end{verbatim}}

\subsection{$K_{18}$, $|\Aut|=4$ --- a representative new class}

A representative of the $144$ classes of P1Fs of $K_{18}$ with automorphism group of
order $4$ (cyclic, $C_4$). To our knowledge no class of this kind has previously been
recorded. (One of the $179$ order-4 classes; the full list is in the repository.)

{\footnotesize\begin{verbatim}
F1  (0,1)(2,3)(4,5)(6,7)(8,9)(10,11)(12,13)(14,15)(16,17)
F2  (0,2)(1,4)(3,6)(5,8)(7,10)(9,12)(11,14)(13,16)(15,17)
F3  (0,3)(1,5)(2,7)(4,9)(6,11)(8,13)(10,15)(12,17)(14,16)
F4  (0,4)(1,3)(2,8)(5,10)(6,9)(7,17)(11,13)(12,14)(15,16)
F5  (0,5)(1,2)(3,9)(4,11)(6,16)(7,8)(10,12)(13,15)(14,17)
F6  (0,6)(1,15)(2,5)(3,17)(4,12)(7,11)(8,14)(9,16)(10,13)
F7  (0,7)(1,14)(2,16)(3,4)(5,13)(6,10)(8,17)(9,15)(11,12)
F8  (0,8)(1,10)(2,14)(3,11)(4,13)(5,16)(6,15)(7,12)(9,17)
F9  (0,9)(1,11)(2,10)(3,15)(4,17)(5,12)(6,13)(7,14)(8,16)
F10 (0,10)(1,9)(2,17)(3,8)(4,15)(5,11)(6,12)(7,16)(13,14)
F11 (0,11)(1,8)(2,9)(3,16)(4,10)(5,14)(6,17)(7,13)(12,15)
F12 (0,12)(1,16)(2,11)(3,5)(4,7)(6,14)(8,15)(9,13)(10,17)
F13 (0,13)(1,17)(2,4)(3,10)(5,6)(7,15)(8,12)(9,14)(11,16)
F14 (0,14)(1,6)(2,15)(3,12)(4,8)(5,7)(9,11)(10,16)(13,17)
F15 (0,15)(1,7)(2,13)(3,14)(4,6)(5,9)(8,10)(11,17)(12,16)
F16 (0,16)(1,13)(2,12)(3,7)(4,14)(5,17)(6,8)(9,10)(11,15)
F17 (0,17)(1,12)(2,6)(3,13)(4,16)(5,15)(7,9)(8,11)(10,14)
\end{verbatim}}

\subsection{$K_{20}$, $|\Aut|=342$ --- the $\AGL(1,19)$ factorisation}

The algebraic P1F of $K_{20}=K_{19+1}$; automorphism group $x\mapsto ax+b$ over
$\GF(19)$, order $19\cdot18=342$. Recovered among the order-19 classes. Note
$F_{10}$ is the ``reversal'' factor $(0,10)(1,19)(2,18)\cdots(9,11)$ characteristic of
this construction.

{\scriptsize\begin{verbatim}
F1  (0,1)(2,3)(4,5)(6,7)(8,9)(10,11)(12,13)(14,15)(16,17)(18,19)
F2  (0,2)(1,4)(3,6)(5,8)(7,10)(9,12)(11,14)(13,16)(15,18)(17,19)
F3  (0,3)(1,5)(2,7)(4,9)(6,11)(8,13)(10,15)(12,17)(14,19)(16,18)
F4  (0,4)(1,8)(2,6)(3,10)(5,12)(7,14)(9,16)(11,18)(13,19)(15,17)
F5  (0,5)(1,9)(2,11)(3,7)(4,13)(6,15)(8,17)(10,19)(12,18)(14,16)
F6  (0,6)(1,12)(2,10)(3,14)(4,8)(5,16)(7,18)(9,19)(11,17)(13,15)
F7  (0,7)(1,13)(2,15)(3,11)(4,17)(5,9)(6,19)(8,18)(10,16)(12,14)
F8  (0,8)(1,16)(2,14)(3,18)(4,12)(5,19)(6,10)(7,17)(9,15)(11,13)
F9  (0,9)(1,17)(2,19)(3,15)(4,18)(5,13)(6,16)(7,11)(8,14)(10,12)
F10 (0,10)(1,19)(2,18)(3,17)(4,16)(5,15)(6,14)(7,13)(8,12)(9,11)
F11 (0,11)(1,18)(2,16)(3,19)(4,14)(5,17)(6,12)(7,15)(8,10)(9,13)
F12 (0,12)(1,15)(2,17)(3,13)(4,19)(5,11)(6,18)(7,9)(8,16)(10,14)
F13 (0,13)(1,14)(2,12)(3,16)(4,10)(5,18)(6,8)(7,19)(9,17)(11,15)
F14 (0,14)(1,11)(2,13)(3,9)(4,15)(5,7)(6,17)(8,19)(10,18)(12,16)
F15 (0,15)(1,10)(2,8)(3,12)(4,6)(5,14)(7,16)(9,18)(11,19)(13,17)
F16 (0,16)(1,7)(2,9)(3,5)(4,11)(6,13)(8,15)(10,17)(12,19)(14,18)
F17 (0,17)(1,6)(2,4)(3,8)(5,10)(7,12)(9,14)(11,16)(13,18)(15,19)
F18 (0,18)(1,3)(2,5)(4,7)(6,9)(8,11)(10,13)(12,15)(14,17)(16,19)
F19 (0,19)(1,2)(3,4)(5,6)(7,8)(9,10)(11,12)(13,14)(15,16)(17,18)
\end{verbatim}}

\subsection{$K_{18}$, $|\Aut|=17$ --- the cyclic (Skolem-starter) factorisation}

The second of the two 1-rotational P1Fs of $K_{18}$ (\S5.5); its automorphism group is the
cyclic group $\mathbb{Z}_{17}$ generated by $x\mapsto x+1\pmod{17}$ (vertex $17$ playing
$\infty$, fixed). It is the cyclic development of the starter obtained from the order-$8$ Skolem
sequence $(1,1,3,7,8,3,2,6,2,5,7,4,8,6,5,4)$ \cite{PikeShalaby}, and our method recovers it among
the order-$17$ classes (the other being the $\AGL(1,17)$ member of Appendix~A.1). The
canonical form below is verbatim solver output; like A.1 its first two factors are the canonical
$F_1,F_2$, but it is a distinct isomorphism class.

{\footnotesize\begin{verbatim}
F1  (0,1)(2,3)(4,5)(6,7)(8,9)(10,11)(12,13)(14,15)(16,17)
F2  (0,2)(1,4)(3,6)(5,8)(7,10)(9,12)(11,14)(13,16)(15,17)
F3  (0,3)(1,5)(2,8)(4,11)(6,15)(7,13)(9,17)(10,12)(14,16)
F4  (0,4)(1,9)(2,16)(3,7)(5,15)(6,14)(8,10)(11,12)(13,17)
F5  (0,5)(1,10)(2,14)(3,16)(4,9)(6,12)(7,15)(8,17)(11,13)
F6  (0,6)(1,16)(2,13)(3,5)(4,10)(7,9)(8,15)(11,17)(12,14)
F7  (0,7)(1,13)(2,4)(3,12)(5,17)(6,11)(8,14)(9,16)(10,15)
F8  (0,8)(1,6)(2,15)(3,13)(4,12)(5,10)(7,16)(9,11)(14,17)
F9  (0,9)(1,12)(2,6)(3,10)(4,8)(5,11)(7,17)(13,14)(15,16)
F10 (0,10)(1,7)(2,9)(3,14)(4,13)(5,16)(6,17)(8,11)(12,15)
F11 (0,11)(1,17)(2,12)(3,15)(4,7)(5,14)(6,9)(8,16)(10,13)
F12 (0,12)(1,15)(2,5)(3,8)(4,16)(6,13)(7,11)(9,14)(10,17)
F13 (0,13)(1,2)(3,17)(4,14)(5,9)(6,8)(7,12)(10,16)(11,15)
F14 (0,14)(1,11)(2,17)(3,4)(5,7)(6,10)(8,13)(9,15)(12,16)
F15 (0,15)(1,8)(2,10)(3,9)(4,6)(5,13)(7,14)(11,16)(12,17)
F16 (0,16)(1,14)(2,7)(3,11)(4,17)(5,6)(8,12)(9,10)(13,15)
F17 (0,17)(1,3)(2,11)(4,15)(5,12)(6,16)(7,8)(9,13)(10,14)
\end{verbatim}}

\subsection{$K_{18}$, $|\Aut|=16$ --- a P1F with a fixed-point-free $2^9$ involution}

The example settling the last open order-$2$ case (\S5.4): a perfect $1$-factorisation of
$K_{18}$ whose automorphism group (of order $16$) contains the \textbf{fixed-point-free
involution}
\[
\alpha=(0\,1)(2\,4)(3\,5)(6\,8)(7\,9)(10\,12)(11\,13)(14\,16)(15\,17),
\]
of cycle type $2^9$. It was recovered within the order-$4$ enumeration (since $4\mid 16$),
independently reproduced by a separate run as a byte-identical canonical matrix, and
verified directly: all $\binom{17}{2}=136$ factor-pairs are Hamiltonian, and $\alpha$
permutes the factors below as \textbf{$9$ fixed factors and $4$ swapped pairs} --- the structure
forced for any $2^9$-invariant P1F. The group of order $16$ in fact contains two such $2^9$
involutions (and one of type $2^8 1^2$); $\alpha$ above is one of them.

{\footnotesize\begin{verbatim}
F1  (0,1)(2,3)(4,5)(6,7)(8,9)(10,11)(12,13)(14,15)(16,17)
F2  (0,2)(1,4)(3,6)(5,8)(7,10)(9,12)(11,14)(13,16)(15,17)
F3  (0,3)(1,5)(2,7)(4,9)(6,11)(8,13)(10,15)(12,17)(14,16)
F4  (0,4)(1,3)(2,12)(5,11)(6,14)(7,8)(9,17)(10,16)(13,15)
F5  (0,5)(1,2)(3,13)(4,10)(6,9)(7,15)(8,16)(11,17)(12,14)
F6  (0,6)(1,9)(2,16)(3,4)(5,15)(7,12)(8,11)(10,17)(13,14)
F7  (0,7)(1,8)(2,5)(3,17)(4,14)(6,13)(9,10)(11,16)(12,15)
F8  (0,8)(1,6)(2,11)(3,15)(4,13)(5,17)(7,16)(9,14)(10,12)
F9  (0,9)(1,7)(2,14)(3,10)(4,16)(5,12)(6,17)(8,15)(11,13)
F10 (0,10)(1,13)(2,9)(3,12)(4,11)(5,6)(7,14)(8,17)(15,16)
F11 (0,11)(1,12)(2,13)(3,8)(4,7)(5,10)(6,15)(9,16)(14,17)
F12 (0,12)(1,10)(2,6)(3,14)(4,8)(5,16)(7,9)(11,15)(13,17)
F13 (0,13)(1,11)(2,15)(3,7)(4,17)(5,9)(6,8)(10,14)(12,16)
F14 (0,14)(1,16)(2,4)(3,11)(5,13)(6,10)(7,17)(8,12)(9,15)
F15 (0,15)(1,17)(2,10)(3,5)(4,12)(6,16)(7,11)(8,14)(9,13)
F16 (0,16)(1,15)(2,17)(3,9)(4,6)(5,14)(7,13)(8,10)(11,12)
F17 (0,17)(1,14)(2,8)(3,16)(4,15)(5,7)(6,12)(9,11)(10,13)
\end{verbatim}}

Under $\alpha$ the nine fixed factors are
$F_1,F_2,F_3,F_8,F_9,F_{12},F_{13},F_{14},F_{15}$, and the remaining eight form the four
swapped pairs $\{F_4,F_5\},\{F_6,F_7\},\{F_{10},F_{11}\},\{F_{16},F_{17}\}$; the nine
$\alpha$-fixed factors restrict to a perfect $1$-factorisation of $K_{10}$ on the block
quotient, as predicted in \S5.4.

\subsection{$K_{18}$ --- an atomic Latin square of order $17$ (from the $|\Aut|=16$ class)}

The order-$17$ Latin square $I(\mathcal F,0)$ obtained (\S5.6) from the $|\Aut|=16$ perfect
$1$-factorisation of $K_{18}$ by deleting vertex $0$. It is symmetric, idempotent, and
\textbf{atomic} --- row-, column- and symbol-Hamiltonian --- and lies in a main class
distinct from that of the $\AGL(1,17)$ square. Rows, columns and symbols run over
$\{1,\dots,17\}$. Since the source factorisation has $|\Aut|=16$ with $17\nmid 16$, this
square is not diagonally cyclic. Its autoparatopism group has order $\nu=96$ (the
$\AGL(1,17)$ square has $\nu=27744$); this square is paratopic to the order-$17$ square
obtained by folding the second even-starter factorisation $E_B$ of
\cite{BryantMaenhautWanless}, confirming the two are the same main class.

{\footnotesize\begin{verbatim}
  1 10 11  2  3 17 16  4  5  9  8  6  7 12 13 14 15
 10  2  1 17 11 12  3 14  6 15  9 13  4  7 16  8  5
 11  1  3 10 16  2 13  7 15  8 14  5 12 17  6  4  9
  2 17 10  4  1  9 11  6  3 12 16 14  5 15  8 13  7
  3 11 16  1  5 10  8  2  7 17 13  4 15  9 14  6 12
 17 12  2  9 10  6  1 13 14  7  3  8 16  5 11 15  4
 16  3 13 11  8  1  7 15 12  2  6 17  9 10  4  5 14
  4 14  7  6  2 13 15  8  1  5 12 10  3 16  9 17 11
  5  6 15  3  7 14 12  1  9 13  4  2 11  8 17 10 16
  9 15  8 12 17  7  2  5 13 10  1 16 14  4  3 11  6
  8  9 14 16 13  3  6 12  4  1 11 15 17  2  5  7 10
  6 13  5 14  4  8 17 10  2 16 15 12  1 11  7  9  3
  7  4 12  5 15 16  9  3 11 14 17  1 13  6 10  2  8
 12  7 17 15  9  5 10 16  8  4  2 11  6 14  1  3 13
 13 16  6  8 14 11  4  9 17  3  5  7 10  1 15 12  2
 14  8  4 13  6 15  5 17 10 11  7  9  2  3 12 16  1
 15  5  9  7 12  4 14 11 16  6 10  3  8 13  2  1 17
\end{verbatim}}

\end{document}
